\section*{अध्याय \ref{ch:b3:advanced-geochronology} --- उन्नत भूकालमिति}

\begin{solution}{exo:b3:advanced-geochronology:1}
$t_{206} = \ln(1.04500)/\lambda_{238} = \qty{283.8}{\Ma}$; $t_{207} =
\ln(1.3300)/\lambda_{235} = \qty{289.6}{\Ma}$। दोनों आयुओं में अंतर
\qty{6}{\Ma} है: ज़िरकॉन विसंगत है। इसकी $t_{207}$ आयु बड़ी है, इसलिए यह
\omterm{def:b3:advanced-geochronology:concordia}{कॉन्कॉर्डिया} के नीचे आलेखित होता है, सीसा-हानि की ओर।
\end{solution}

\begin{solution}{exo:b3:advanced-geochronology:2}
$J = (e^{5.543\times10^{-10}\times2.8201\times10^7} - 1)/1.3129 =
0.015754/1.3129 = 0.01200$। फिर $t = \ln(1 + 0.01200\times10.000)/\lambda =
\ln 1.12/(5.543\times10^{-10}) = \qty{204.5}{\Ma}$।
\end{solution}

\begin{solution}{exo:b3:advanced-geochronology:3}
$T_{\mathrm{CHUR}} = \ln\bigl(1 + (0.51180 - 0.512638)/(0.1200 -
0.1966)\bigr)/\lambda = \ln(1.010940)/\lambda = \qty{1.67}{\Ga}$;
$T_{\mathrm{DM}} = \ln\bigl(1 + (0.51180 - 0.51315)/(0.1200 - 0.2137)\bigr)/\lambda
= \ln(1.014408)/\lambda = \qty{2.19}{\Ga}$।
\end{solution}

\begin{solution}{exo:b3:advanced-geochronology:4}
क्षय के बिना, $t = N/P_0 = \num{2.0e5}/4.96 = \qty{40.3}{kyr}$। क्षय के साथ,
$t = -\ln(1 - N\lambda/P_0)/\lambda = -\ln(1 - 0.02014)/\lambda =
\qty{40.7}{kyr}$, \qty{1}{\percent} अधिक पुरानी।
\end{solution}

\begin{solution}{exo:b3:advanced-geochronology:5}
छोटे $t$ के लिए, $x \approx \lambda_{235}t$ और $y \approx \lambda_{238}t$, इसलिए
$y \approx (\lambda_{238}/\lambda_{235})x$, ढलान $0.1575$। यदि सीसा
बहुत हाल में खोया गया हो, $t_2 \approx 0$ और \omterm{def:b3:advanced-geochronology:concordia}{निचला अंतःखंड} $C(0)$,
यानी मूल बिंदु। मूल बिंदु से होकर जाने वाली रेखा पर $y/x$ स्थिर है, और
$^{207}\mathrm{Pb}^*/^{206}\mathrm{Pb}^* = x/(137.8\,y)$ भी स्थिर है: हर क्रिस्टल
\omterm{def:b3:advanced-geochronology:concordia}{ऊपरी अंतःखंड} की सीसा--सीसा आयु देता है, यानी संरचना की आयु।
\end{solution}

\begin{solution}{exo:b3:advanced-geochronology:6}
भार $1.5625$, $1$, $1.2346$, $2.0408$, $W = 5.838$; $\bar x =
\qty{296.09}{\Ma}$, $\sigma = 1/\sqrt W = \qty{0.41}{\Ma}$। भारित
वर्ग विचलनों का योग $2.32$ है; $f = 3$ पर, $\mathrm{MSWD} = 0.77$ है, जो
$1 + 2\sqrt{2/3} = 2.63$ से नीचे है: आयुएँ मेल खाती हैं।
\end{solution}

\begin{solution}{exo:b3:advanced-geochronology:7}
$\lambda\zeta g\rho_d\rho_s/\rho_i = 1.55125\times10^{-10}\times350\times0.5
\times\num{1.10e6}\times0.5 = 0.014931$; $t = \ln(1.014931)/\lambda =
\qty{95.5}{\Ma}$। यदि चट्टान तेज़ी से ठंडी हुई, तो यह
ऐपेटाइट के आंशिक अनीलन क्षेत्र से, लगभग 60 से \qty{120}{\celsius} तक, गुज़रने का काल बताती है; यदि
धीरे ठंडी हुई, तो यह मिश्रित आयु है, और इसका पता पथ-लंबाइयों से लगना चाहिए।
\end{solution}

\begin{solution}{exo:b3:advanced-geochronology:8}
प्रति ग्राम: यूरेनियम $20\times10^{-6}/238.05\times N_A = \num{5.06e16}$ परमाणु, जिनमें
\num{5.02e16} यूरेनियम-238 के और \num{3.64e14} यूरेनियम-235 के; थोरियम
\num{1.038e17}, $\lambda_{232} = \ln 2/\num{1.40e10} = \qty{4.95e-11}{yr^{-1}}$;
हीलियम $10^{-9}N_A = \num{6.02e14}$। \omterm{def:b3:advanced-geochronology:cosmogenic}{उत्पादन दर}
$8\lambda_{238}{}^{238}\mathrm U + 7\lambda_{235}{}^{235}\mathrm U +
6\lambda_{232}\mathrm{Th} = (6.23 + 0.25 + 3.08)\times10^7 =
\num{9.57e7}\text{ परमाणु}\,\unit{g^{-1}.yr^{-1}}$ है, और $t = \num{6.02e14}/\num{9.57e7} =
\qty{6.29}{Myr}$। संशोधित: $6.29/0.75 = \qty{8.39}{\Ma}$।
\end{solution}

\begin{solution}{exo:b3:advanced-geochronology:9}
$P_0/\lambda = 4.96/\num{4.994e-7} = \num{9.93e6}\text{ परमाणु}\,\unit{g^{-1}}$, जो
\qty{90}{\percent} तक $\ln 10/\lambda = \qty{4.6}{Myr}$ के बाद पहुँचता है।
\qty{10}{m.Myr^{-1}} पर $= \qty{1e-3}{cm.yr^{-1}}$: $\rho\varepsilon/\Lambda =
2.7\times10^{-3}/160 = \qty{1.69e-5}{yr^{-1}}$, $\mu = \qty{1.74e-5}{yr^{-1}}$
और $N = 4.96/\num{1.74e-5} = \num{2.85e5}\text{ परमाणु}\,\unit{g^{-1}}$, 35 गुना कम।
\end{solution}

\begin{solution}{exo:b3:advanced-geochronology:10}
$W = w_1 + w_2$ लेने पर, $x_1 - \bar x = w_2(x_1 - x_2)/W$ और $x_2 - \bar x =
-w_1(x_1 - x_2)/W$, इसलिए $\sum w_i(x_i - \bar x)^2 = (x_1 - x_2)^2(w_1w_2^2 +
w_2w_1^2)/W^2 = (x_1 - x_2)^2w_1w_2/W = (x_1 - x_2)^2/(\sigma_1^2 +
\sigma_2^2)$; $f = 1$ पर यह MSWD है। यदि दोनों एक ही आयु मापें,
$x_1 - x_2$ का माध्य 0 और प्रसरण $\sigma_1^2 + \sigma_2^2$ है, इसलिए
इसके वर्ग का, उस प्रसरण से भाग देने पर, प्रत्याशित मान 1 होता है।
\end{solution}

\begin{solution}{exo:b3:advanced-geochronology:11}
$C(2700) = (13.283, 0.52019)$, $C(500) = (0.63628, 0.08065)$। $f = (0.52019 -
0.300)/(0.52019 - 0.08065) = 0.501$, इसके सीसे का आधा। $x = 0.499\times13.283 +
0.501\times0.63628 = 6.948$। $t_{206} = \qty{1691}{\Ma}$, $t_{207} =
\qty{2105}{\Ma}$; अनुपात $x/(137.8\,y) = 0.1681$ लेकर पुनरावृत्ति द्वारा
सीसा--सीसा आयु \qty{2537}{\Ma} प्राप्त होती है। इन तीनों का कोई अर्थ नहीं है: ये
दोनों घटनाओं के मिश्रण हैं, और केवल अंतःखंड ही उनकी आयु बताते हैं। सीसा--सीसा
आयु क्रिस्टलन के लिए एक न्यूनतम मान है।
\end{solution}

\begin{solution}{exo:b3:advanced-geochronology:12}
अपरदन न होने पर: $t = -\ln(1 - \num{1.0e6}\times\num{4.994e-7}/4.96)/\lambda =
\qty{212}{kyr}$। स्थायी अवस्था: $\rho\varepsilon/\Lambda = P_0/N - \lambda =
\qty{4.46e-6}{yr^{-1}}$, $\varepsilon = \qty{2.64e-4}{cm.yr^{-1}} =
\qty{2.6}{m.Myr^{-1}}$। इससे कम किसी भी $\varepsilon$ के लिए, $N =
P_0(1 - e^{-\mu t})/\mu$ एक $t$ तय करता है, जो \qty{212}{kyr} से
ऊपर अनंत तक कुछ भी हो सकता है: एक माप में दो अज्ञात राशियाँ होती हैं। किसी अन्य
अर्ध-आयु वाला दूसरा नाभिक, जैसे उसी क्वार्ट्ज़ में एल्युमिनियम-26, या
किसी क्रोड में गहराई पर सांद्रता इन्हें अलग करती है।
\end{solution}

\begin{solution}{pb:b3:advanced-geochronology:1}
\textbf{1.} $t_{206}$ / $t_{207}$, \unit{\Ma} में: Z1 301.9 / 302.2, Z2 284.1 /
286.5, Z3 265.7 / 268.5, Z4 246.1 / 250.1, Z5 225.1 / 229.4। $y$ पर त्रुटि
\qty{0.2}{\percent} की लगभग \qty{0.6}{\Ma} है: Z1 त्रुटि के भीतर संगत है,
बाकी विसंगत हैं, और Z2 से Z5 तक उत्तरोत्तर अधिक।

\textbf{2.} एक बाद की घटना में ज़िरकॉनों ने अपने सीसे का बढ़ते अंश
खो दिया: वे जीवा पर अधिकाधिक नीचे स्थित हैं। जीवा
अवतल \omterm{def:b3:advanced-geochronology:concordia}{कॉन्कॉर्डिया} के नीचे स्थित है, जहाँ $x$, दिए गए $y$ के लिए, बड़ा होता है: $t_{207} > t_{206}$।

\textbf{3.} Z1: $X = 1/0.04794 = 20.859$, $Y = 0.3467/(137.8\times0.04794) =
0.05248$। Z5: $X = 28.145$, $Y = 0.05178$।

\textbf{4.} भार $1.000$, $1.235$, $1.563$, $1.563$, $2.041$ (
$10^8$ में); $\bar x = 0.29402$, $\bar y = 0.040892$, $S_{xx} = \num{7.845e5}$।

\textbf{5.} $b = 0.13260$, $a = \bar y - b\bar x = 0.001905$।

\textbf{6.} \qty{310}{\Ma} पर, $\mathrm dy/\mathrm dx = 0.1575\,
e^{(\lambda_{238} - \lambda_{235})t} = 0.1218$। \omterm{def:b3:advanced-geochronology:concordia}{डिस्कॉर्डिया} केवल
\qty{9}{\percent} \omterm{def:b3:advanced-geochronology:concordia}{कॉन्कॉर्डिया} से अधिक तीखी है: वह उसे छोटे कोण पर काटती है,
और \omterm{def:b3:advanced-geochronology:concordia}{ऊपरी अंतःखंड} सटीक नहीं होगा।

\textbf{7.} \qty{310}{\Ma} से आरंभ करके न्यूटन विधि: $t = \qty{311.2}{\Ma}$।

\textbf{8.} \qty{100}{\Ma} से आरंभ करके: $t = \qty{99.7}{\Ma}$।

\textbf{9.} $(y_i - a - bx_i)/\sigma_i = 0.64$, $-0.91$, $0.35$, $-0.24$,
$0.17$; $\mathrm{MSWD} = 1.435/3 = 0.48$, $2.63$ की सीमा से नीचे, जो $f = 3$ के लिए है:
ज़िरकॉन एक ही रेखा पर स्थित हैं।

\textbf{10.} $s_b = 1/\sqrt{\num{7.845e5}} = 0.00113$।

\textbf{11.} $b - 2s_b = 0.13034$ लेने पर \omterm{def:b3:advanced-geochronology:concordia}{ऊपरी अंतःखंड}
\qty{322.9}{\Ma} है; $b + 2s_b = 0.13485$ लेने पर, \qty{302.6}{\Ma}। आयु
$311 \pm \qty{10}{\Ma}$ है, दो मानक विचलन पर।

\textbf{12.} $f = (x(311.2) - x_i)/(x(311.2) - x(99.7))$: Z1 \qty{4.7}{\percent},
Z5 \qty{41}{\percent}।

\textbf{13.} $J = 0.015754/1.3129 = 0.0119999$, $0.01200$।

\textbf{14.} चरण 1: \qty{150.0}{\Ma}; चरण 4: \qty{296.6}{\Ma}; चरण 10:
\qty{295.0}{\Ma}।

\textbf{15.} चरण 4 से 10: लगातार सात चरण, जिनमें \qty{82}{\percent}
आर्गन-39 है, कोई प्रवृत्ति नहीं। चरण 1 से 3 सीढ़ी के समान बढ़ते हैं।

\textbf{16.} समान त्रुटियाँ: माध्य $2072.9/7 = \qty{296.1}{\Ma}$, $\sigma =
1.1/\sqrt7 = \qty{0.42}{\Ma}$; वर्गित विचलनों का योग $3.83$ है,
$\mathrm{MSWD} = 3.83/(1.21\times6) = 0.53$, $2.15$ से नीचे।

\textbf{17.} $\bar R = \sum f_iR_i = 14.270$, $t = \qty{285.2}{\Ma}$:
किनारों से खोया आर्गन, जो पहले चरणों में मुक्त हुआ, माध्य को घटाता है,
जैसे वह पोटैशियम--आर्गन आयु को घटाता।

\textbf{18.} $(700 - 375)/(311.2 - 296.1) = \qty{21.5}{\celsius.Myr^{-1}}$
(व्यापक परास के साथ, क्योंकि ज़िरकॉन आयु $\pm\qty{10}{\Ma}$ तक ज्ञात है)।

\textbf{19.} \qty{100}{\Ma} के निकट ग्रेनाइट अपने मेसोज़ोइक
आवरण के नीचे दब गया था, और किसी तापन या द्रवों के परिसंचरण से ज़िरकॉनों ने सीसा खो दिया।
बायोटाइट ने अपने किनारों से आर्गन खोया, जिसे सीढ़ी-आकार दर्शाता है; इसका
सबसे निचला चरण, \qty{150}{\Ma}, पुराने और पुनः सेट आर्गन का मिश्रण है। पठार
बच गया: बायोटाइट इतना नहीं तपा कि अपना सारा आर्गन खो दे।

\textbf{20.} आयु के विरुद्ध ऊँचाई की न्यूनतम वर्ग ढलान
\qty{20.1}{m.Myr^{-1}} है।

\textbf{21.} $(65 - 10)/30 = \qty{1.83}{km}$; \qty{43.1}{\Ma} से, माध्य
$1830/43.1 = \qty{42}{m.Myr^{-1}}$ है।

\textbf{22.} $\rho\varepsilon/\Lambda = P_0/N - \lambda = \num{3.415e-5} -
\num{0.050e-5} = \qty{3.365e-5}{yr^{-1}}$; $\varepsilon = \num{3.365e-5}\times
160/2.7 = \qty{1.99e-3}{cm.yr^{-1}} = \qty{19.9}{m.Myr^{-1}}$।

\textbf{23.} $P_0\Lambda/(\rho N) = \qty{20.2}{m.Myr^{-1}}$, \qty{1.5}{\percent}
अधिक: इस दर पर क्षय का लगभग कोई प्रभाव नहीं पड़ता।

\textbf{24.} $\Lambda/(\rho\varepsilon) = 160/(2.7\times\num{1.99e-3}) =
\qty{30}{kyr}$। ग्रेनाइट लगभग \qty{20}{m.Myr^{-1}} की दर से 52 और
\qty{43}{\Ma} के बीच ऊपर उठा, और तब से औसतन \qty{42}{m.Myr^{-1}} पर, तथा अब
\qty{20}{m.Myr^{-1}} पर अपरदित हो रहा है: सेनोज़ोइक में किसी समय
आवरण-हटाव अधिक तेज़ था।

\textbf{25.} वार्डन ग्रेनाइट स्फटिकीकृत हुआ
\textbf{\qty[text-series-to-math=true, separate-uncertainty=true]{311 \pm 10}{\Ma}} पहले (\omterm{def:b3:advanced-geochronology:concordia}{ऊपरी अंतःखंड},
दो मानक विचलन पर)।
\end{solution}
